\section{Supporting Information}
Details of practical method and data analysis which cramped the
 presentation of material in the present paper have been gathered
 and elucidated within this section.

\subsection{Notes on Synthesis of Diitaconate}
\label{subsec:NoOnSyOfDi}
Dimethyl itaconate was synthesised by esterification of itaconic acid
 with methanol by the Fischer method, here we describe the practical
 method by which this was achieved.
Figure~\ref{fig:AppendixSynthesisRigDMI_Lettered} shows how this reaction mixture
 was placed in the lower round-bottomed flask
 (Fig.~\ref{fig:AppendixSynthesisRigDMI_Lettered}D), and brought to a  vigorous boil
 at around 130\unit{\degreeCelsius}.
Vapours of both water and methanol passed up an insulated glass tube
 (Fig.~\ref{fig:AppendixSynthesisRigDMI_Lettered}E) and were captured in a 
 distillation column (Fig.~\ref{fig:AppendixSynthesisRigDMI_Lettered}G) which was 
 connected to a reboiler (Fig.~\ref{fig:AppendixSynthesisRigDMI_Lettered}F) at its 
 base.
The temperature at the top of the column (Fig.~\ref{fig:AppendixSynthesisRigDMI_Lettered}H)
 was maintained at around 70\unit{\degreeCelsius} by adjusting the 
 temperature of the reboiler (120\unit{\degreeCelsius}) and the level of 
 insulation applied around the column; this ensured a steady stream of 
 methanol-rich vapours passing through the Claisen head 
 (Fig.~\ref{fig:AppendixSynthesisRigDMI_Lettered}I) and re-condensing within the 
 Liebig condenser (Fig.~\ref{fig:AppendixSynthesisRigDMI_Lettered}J).
Methanol-rich liquid was therewith returned to the lower round-bottomed
 flask via a dropping funnel (Fig.~\ref{fig:AppendixSynthesisRigDMI_Lettered}K),
 whilst water was caught in the reboiler
 (Fig.~\ref{fig:AppendixSynthesisRigDMI_Lettered}F), pushing the reversible
 esterification reaction toward the desired diester product.

\noindent
\begin{minipage}{0.485\textwidth}
    \includegraphics[width=\textwidth]{SynthesisRigDMI_Lettered.png}
    \captionsetup{format=custom,width=\textwidth,skip=0pt}
    \captionof{figure}{
        \footnotesize{
            Apparatus for synthesis of dimethyl itaconate.
            (A) Hotplate stirrer
            (B) Aluminium bowl
            (C) Temperature controller probe
            (D) Three-necked round-bottomed flask containing the reaction
                mixture
            (E) Insulated glass tube
            (F) Single-necked round-bottomed flask acting as a distillation
                column reboiler containing water
            (G) Distillation column
            (H) Thermometer
            (I) Four-way Claisen head
            (J) Liebig condenser
            (K) Dropping funnel for recycled methanol
        }
    }
    \label{fig:AppendixSynthesisRigDMI_Lettered}
\end{minipage}

Before initiating the reaction by starting both the hotplate stirrers
 (Fig.~\ref{fig:AppendixSynthesisRigDMI_Lettered}A), an argon-rich atmosphere was
 generated by using a Schlenk line to cycle a vacuum (500~\unit{\milli\bar})
 and flow of argon gas (350~\unit{\milli\litre\per\minute}); this was to 
 discourage unwanted oxidation reactions and reduce the chances of vapour 
 ignition.
During the reaction however a valve was opened to allow pressure 
 equilibration between the reaction vessels and the atmosphere in order to
 mitigate the risk of explosion.
Methanol cycling was halted 3.5~\unit{\hour} into the reaction by closing 
 the dropping funnel (Fig.~\ref{fig:AppendixSynthesisRigDMI_Lettered}K), after
 which all heating was ceased on reaching 4~\unit{\hour}.

Having trialled several solvents (e.g. cyclohexane, n-hexane), toluene 
 was found to be the most suitable as its relatively higher affinity for 
 dimethyl itaconate during liquid-liquid extraction led to acceptable
 yields at conclusion of workup.
The final product mixture retrieved from the round-bottomed flask
 (Fig.~\ref{fig:AppendixSynthesisRigDMI_Lettered}D) was placed in a separating funnel
 with toluene in a 1:3 volume ratio.
The separating funnel contents were washed with water three times, saturated
 sodium bicarbonate solution twice, and saturated sodium chloride once;
 each washing step involved a volume ratio of 2:1 between the previously
 obtained organic fraction and the next volume of washing liquid used.
A gravity filtration step was then carried out between the water and sodium
 bicarbonate washing steps to remove insolubles.
Following washing, the organic fraction was removed to an Erlenmeyer flask
 and dried using anhydrous magnesium sulphate; thereafter removed using
 gravity filtration.
Dry organics were placed in a rotary evaporator and heated at
 40\unit{\degreeCelsius} whilst the pressure was reduced from an initial
 70~\unit{\milli\bar} to 20~\unit{\milli\bar} over an average period of
 8~\unit{\hour}; this removed all but a small mass of toluene and water
 from the liquid product.
From this product, dimethyl itaconate was crystallised out at room
 temperature by addition of a nucleating agent.
The nucleating agent was crystalline dimethyl itaconate obtained directly
 from a particularly pure batch of rotary evaporation product which had
 crystallised in a fridge at 3\unit{\degreeCelsius}.
The above method afforded a 45\unit{\percent} yield of product,although 
 higher values may be possible if an extended period of boiling were 
 adopted.

\subsection{Notes on the TLC of Diitaconates}
\label{subsec:NoOnThTLCOfDii}
Initial estimates of purity were garnered from spotting chemicals onto
 thin layer chromatography plates (polyester with a layer of silica gel and
 ultraviolet indicator; 6~\unit{\milli\metre} pore diameter and
 0.2~\unit{\milli\metre} layer thickness) and using a mixture of
 dichloromethane and methanol (2.5:7.5 v/v respectively) as the mobile phase.
On completion, spot(s) could be observed under ultraviolet light
 (\(\lambda=\)254~\unit{\nano\metre}) and their retardation factors
 (\(R_{\mathrm{F}}\)) calculated using equation \ref{eq:RetardationFactor},
 relating this metric with the distance migrated by the substance in
 question (\(b\) in \unit{\milli\metre}) and that travelled by the
 mobile phase (\(a\) in \unit{\milli\metre}) \autocite{Ettre1993}.
\begin{equation}\label{eq:RetardationFactor}
    R_{\mathrm{F}}=\frac{b}{a}
\end{equation}
Having solvated and diluted pure itaconic acid to a 1.25\unit{\percent}~w/w
 solution using stock mobile phase, this solution was spotted on a plate- this
 produced a single spot with an \(R_{\mathrm{F}}\) value of 0.79.
Meanwhile, having spotted a 5\unit{\percent}~w/w solution of the final
 crystalline product (expected to be dimethyl itaconate), a single spot with
 an \(R_{\mathrm{F}}\) of 0.91 was obtained.
These results provided benchmarks by which reaction progress could be
 gauged as well as supporting the hypothesis that these two crystalline
 materials were reasonably pure, although the
 presence of non-fluorescent or trace chemicals cannot be ruled out
 using TLC alone.

Samples taken periodically from the reaction vessel during the course of
 reflux revealed a change in composition: the first sample
 (retrieved after 15~\unit{\minute} of heating) presented a single spot with
 an \(R_{\mathrm{F}}\) value of 0.81, while a final sample (retrieved after 
 4~\unit{\hour}) presented two spots with \(R_{\mathrm{F}}\) values of 0.90 and 0.81
 respectively.
Given the similarity of these samples with those of the crystalline
 reactant and product analysed previously, these results suggest that
 partial conversion of the reactant itaconic acid into the desired product
 dimethyl itaconate was achieved.

\subsection{Notes on Synthesis of Unsaturated Polyester Resin}
\label{subsec:NoOnSyOfUnPoRe}
Figure~\ref{fig:AppendixSynthesisRigUP_Lettered} shows the apparatus used to
 synthesise the unsaturated polyester.
A Dean-Stark trap (Fig.~\ref{fig:AppendixSynthesisRigUP_Lettered}I) was used to
 allow the `azeotropic method' of polyester synthesis (as opposed to
 the `dry' method) which reduced the time taken to synthesise a
 high molecular weight polyester \autocite{Parkyn1967}.
In this method, an entraining agent (mixed xylenes) which can form an
 azeotrope with the water was added to the mixture (5\unit{\percent} w/w)
 \autocite{Korbar1993,Bodnar1990}.
This increases the rate at which water is condensed and captured by the 
 Dean-Stark trap, driving the reaction forward by removing the unwanted 
 byproduct of water from the reaction mixture.
Furthermore, argon gas was used to sparge the reaction mixture
 (Fig.~\ref{fig:AppendixSynthesisRigUP_Lettered}E) facilitating the removal of
 water from the bulk mixture and mitigating the combustion risk associated
 with high temperature xylenes \autocite{Parkyn1967}.

The synthesis procedure was initiated by charging the four-necked 
 round-bottomed flask (Fig.~\ref{fig:AppendixSynthesisRigUP_Lettered}D) with the 
 requisite masses of the two monomers and xylenes, opening the argon taps 
 to provide a flow of sparging gas at the rate of 5 bubbles per second 
 through the argon injector (Fig.~\ref{fig:AppendixSynthesisRigUP_Lettered}E)
 \autocite{Parkyn1967}, opening the water taps to actively cool the 
 Liebig condenser (Fig.~\ref{fig:AppendixSynthesisRigUP_Lettered}I), setting the 
 overhead stirrer to a speed of 50~rpm 
 (Fig.~\ref{fig:AppendixSynthesisRigUP_Lettered}G) \autocite{Parkyn1967}, and 
 bringing the reaction mixtures temperature 
 (Fig.~\ref{fig:AppendixSynthesisRigUP_Lettered}F) up to 160\unit{\degreeCelsius}
 using the hotplate (Fig.~\ref{fig:AppendixSynthesisRigUP_Lettered}A).
During the procedure the mixture's temperature was raised monotonically 
 by 10\unit{\degreeCelsius\per\hour} until it reached a maximum of
 210\unit{\degreeCelsius}, at which it was held.

\noindent
\begin{minipage}{0.485\textwidth}
    \includegraphics[width=\textwidth]{SynthesisRigUP_Lettered.png}
    \captionsetup{format=custom,width=\textwidth,skip=0pt}
    \captionof{figure}{
        \footnotesize{
            Apparatus for synthesis of unsaturated poly(propyl itaconate).
            (A) Hotplate stirrer
            (B) Aluminium bowl
            (C) Temperature controller probe
            (D) Four-necked round-bottomed flask containing the reaction
                mixture
            (E) Stainless-steel argon injector
            (F) Probes for batch and vapour tamperatures
            (G) Overhead stirrer
            (H) Dean-Stark trap
            (G) Liebig condenser.
        }
    }
    \label{fig:AppendixSynthesisRigUP_Lettered}
\end{minipage}

The reaction progress was ascertained by periodically retrieving a sample of the
 reaction mixture and ascertaining its acid number, in units of
 \unit{\milli\gram\per\gram}, describing the amount of potassium hydroxide
 required to neutralise one gram of sampled unsaturated polyester.
This approximate form of end-group analysis is commonly used in unsaturated
 polyester synthesis \autocite{Pantic2022,Spasojevic2021,Costa2017,Parkyn1967}
 and provides a proxy value estimating the number of un-esterified acid
 groups remaining within the unsaturated polyester.
An American Society for Testing and Materials standard (``ASTM D4662-20'')
 was used for this purpose.
In-line with common practice, heating of the polyesterifying mixture was
 ceased at an acid value of 50 \autocite{Parkyn1967,Bodnar1990,Korbar1993}.
On cessation of heating, an outlet was opened for egress of remnant xylenes
 and the unsaturated polyester was allowed to cool in an air bath.

Since the viscosity of a cooling unsaturated polyester will increase 
 monotonically until it reaches room temperature, where it exists as a
 solid thermoplastic mass, it is common to prepare the unsaturated
 polyester resin immediately after cessation of heating by adding
 reactive diluent directly to the cooling unsaturated polyester
 \autocite{Parkyn1967}.
Unfortunately, reactive diluents are liable to
 polymerise during this addition causing a premature gelation.
Mitigation against hardening into thermoplastic or
 premature gelation into rubber can be achieved by addition of an 
 addition-polymerisation inhibiting agent at a suitable intermediate
 temperature \autocite{Parkyn1967,Ahamad2001,Cherian2004}.
In this work, we awaited a temperature of 100\unit{\degreeCelsius} before 
 adding the below specified amounts of inhibitor and reactive diluent.

Here, we used the commonly-used free-radical scavenger hydroquinone
 as an inhibiting agent, which was added so as to make up
 0.06\unit{\percent}~w/w of the cooling unsaturated polyester mixture.
Hydroquinone was dissolved in a small amount of propylene glycol before
 addition to ensure homogenous mixing into the bulk
 \autocite{Parkyn1967}.
The inhibitor percentage was determined through a process of
 trial-and-error, where lower percentages tended to facilitate premature
 gelation of unsaturated polyester resins on the shelf, whilst higher
 percentages entirely prevented the intended curing procedures for these 
 resins.
The percentage of inhibiting agent was above that which is
 commonly used in the synthesis of conventional unsaturated polyester
 resins (approximately 0.02\unit{\percent}) \autocite{Parkyn1967,Cherian2004}.
This was a consequence of both the relatively higher proportion of
 unsaturated sites along the polyester backbone and a higher copolymerisation
 affinity between the reactive diluent and the sites of unsaturation.

Following introduction and homogenisation of inhibiting agent into
 the unsaturated polyester mixture, the reactive diluent (dimethyl itaconate)
 was added such that it made up 30\unit{\percent}~w/w of the final
 unsaturated polyester resin.

\subsection{Notes on Gel Permeation Chromatography}
\label{subsec:NoOnGePeCh}
In a separate experimental series (not considered in the context of the
 present work), a similar polyester was synthesised with a reduced content
 of itaconate moieties in its backbone.
This was achieved by replacing some of the itaconic acid used during
 synthesis with succinic acid, producing poly(propylene glycol succinate
 itaconate).
Figure~\ref{fig:AppendixAlternativeSampleGPC} depicts the GPC results obtained
 on analysing this alternative formulation.
On comparison with the GPC results obtained for poly(propylene glycol
 itaconate) as detailed in the present work it becomes apparent that
 this alternative formulation exhibits little if any premature crosslinking
 and associated network polymers.
As the closest approximation to the distribution associated a
 poly(propylene glycol itaconate) polyester that didn't suffer 
 network polymer genesis, we re-scaled the tail section of
 Figure~\ref{fig:AppendixAlternativeSampleGPC} and inserted this into the
 apparent distribution for poly(propylene glycol itaconate)
 presented in the present work.
This allowed the calculation of both number and weight average molecular
 weights for the hypothetically uncrosslinked polyester.

\noindent
\begin{minipage}{0.485\textwidth}
    \includegraphics[width=\textwidth]{Sample_PPSI.png}
    \captionsetup{format=custom,width=\textwidth,skip=0pt}
    \captionof{figure}{
        \footnotesize{
            A graph depicting the distribution of mass across various
            polymer chain lengths for an alternative polyester sample
            tested using GPC.
        }
    }
    \label{fig:AppendixAlternativeSampleGPC}
\end{minipage}

An alternative line of investigation undertaken with respect to the
 GPC of these polyesters was that associated with the deconvolution of
 the apparent distribution.
By identification of three physical peaks (Figure~\ref{fig:SampleGPC}A/B/C)
 mathematical deconvolution was undertaken by fitting a log normal
 distribution to the three curves in isolation (by cropping the apparent
 distribution suitably), before fitting the three in unison (in an additive
 fashion) to the apparent distribution.
Figure~\ref{fig:DeconvGPC} shows the three deconvoluted peaks and their fit
 relative to the apparent distribution.

\noindent
\begin{minipage}{0.485\textwidth}
    \includegraphics[width=\textwidth]{Deconvoluted.png}
    \captionsetup{format=custom,width=\textwidth,skip=0pt}
    \captionof{figure}{
        \footnotesize{
            A graph of all three deconvoluted peaks which were
            identified from within the apparent distribution.
        }
    }
    \label{fig:DeconvGPC}
\end{minipage}

Having deconvoluted the three peaks and mused on the possibility of
 cyclisation as the mechanism behind the existence of the two smaller
 peaks, we investigated various cyclisation mechanisms.
Common to all polyesters are the ring closure and back biting
 mechnanisms \autocite{Ahjopalo2000} whilst Ordelt reactions have been
 identified as particularly possible in the case of unsaturated polyesters
 incorporating itaconate moieties \autocite{Molina2020}.
Since back biting and ring closure both result in the same structure we
 only modelled the latter.
Figure~\ref{fig:RCCM} visualises ring closure: Fig~\ref{fig:RCCM}A
 shows a linear section of polymer with a start length \(L_{0}\) of 1,
 Fig~\ref{fig:RCCM}B shows a perfectly circular loop with a maximum
 length \(L_{1}\) of 0.32, whilst Fig~\ref{fig:RCCM}C considers the
 extreme scenario in which the loop flattens entirely producing a 
 maximum length \(L_{1}\) of 0.5.

\noindent
\begin{minipage}{0.485\textwidth}
    \includegraphics[width=\textwidth]{RingClosureStringModel.png}
    \captionsetup{format=custom,width=\textwidth,skip=0pt}
    \captionof{figure}{
        \footnotesize{
            A schematic describing the model employed when considering
            the ring-closure cyclisation mechanism.
        }
    }
    \label{fig:RCCM}
\end{minipage}

Figure~\ref{fig:ORDM} visualises Ordelt-based cyclisation:
 Fig~\ref{fig:ORDM}A shows the linear section of polymer with an initial
 length of \(L_{0}=1\), whilst an annotation below indicates the random
 nature by which an end hydroxyl group might react with an itaconate
 unsaturation anywhere along the polymer's backbone.
Fig~\ref{fig:ORDM}B describes a perfectly circular loop situation where
 an end group has reacted at position \(L_{1}=0.75\), producing a final
 maximum length \(L_{2}\) of 0.92.
Fig~\ref{fig:ORDM}C describes the extreme scenario for the previous
 hypothetical \(L_{1}=0.75\) situation, where the loop flattens completely
 and generates a final maximum length \(L_{2}=0.88\).
Fig~\ref{fig:ORDM}C1/C2 provide the two extreme scenarios of ring
 flattening at both ends of the possible range of \(L_{1}\) values, providing
 the ultimately highest and lowest final maximum lengths possible under the
 Ordelt reaction mechanism \(L_{2}=[0.99,0.51]\).

\noindent
\begin{minipage}{0.485\textwidth}
    \includegraphics[width=\textwidth]{OrdeltStringModel.png}
    \captionsetup{format=custom,width=\textwidth,skip=0pt}
    \captionof{figure}{
        \footnotesize{
            A schematic describing the model employed when considering
            the Ordelt cyclisation mechanism.
        }
    }
    \label{fig:ORDM}
\end{minipage}

Given the set initial length \(L_{0}=1\), Figure~\ref{fig:HDVM} describes
 the estimation of hydrodynamic volume for each of the scenarios (linear,
 Ordelt cyclised, and ring closure cyclised).
Here, for each of the modelled scenarios we assumed conformations with
 the highest maximum lengths, we also assumed an average length between
 the maximum and minimum situations for the Ordelt scenario, whilst
 proxied hydrodynamic volumes were calculated by assuming the swept 
 spherical area of conformations.
A simple comparison of the final hydrodynamic volume proxy \(V_{h}\) 
 values shows a clear hierarchy between the linear (\(V_{h}=1.05\)),
 the Ordelt (\(V_{h}=0.59\)), and the ring closure (\(V_{h}=0.26\)) 
 scenarios.
Since the volume scales as a power of the dimensions, small changes
 in the total final length of the conformation could possible cause
 drastic changes in the hydrodynamic volume.
Figure~\ref{fig:HDVM} shows how a polyester molecule of the same
 molecular weight (modelled here as length) can exhibit a variety of 
 hydrodynamic volumes due to both cyclisation mechanisms and coformations 
 adopted thereafter.

\noindent
\begin{minipage}{0.485\textwidth}
    \includegraphics[width=\textwidth]{VolumeRodSphereModel.png}
    \captionsetup{format=custom,width=\textwidth,skip=0pt}
    \captionof{figure}{
        \footnotesize{
            A schematic describing the simple model by which
            hydrodynamic volume of polymers was estimated.
        }
    }
    \label{fig:HDVM}
\end{minipage}


\subsection{Notes on Mechanical Test Data Processing}
Here is an unabridged overview of all the data processing undertaken
 during the mechanical testing of samples.

\paragraph{Euler Buckling}
To mitigate against concerns of Euler buckling taking place in samples of
 the proposed dimensions, we calculated the critical load based on
 Euler's critical load equation and some estimated properties for
 conventional polyester thermosets \autocite{Euler1744}.
Equation \ref{eq:AppendixEulerCriticalLoad} shows how the critical load 
 (\(P_{cr}\) in \unit{\newton}) is related to the elastic modulus of the
 material (\(E\) in \unit{\pascal}), the second moment of area (\(I\) in
 \unit{\metre^{4}}),
 the effective length factor (\(K\), unitless), and the initial length of
 the column in question (\(L\) in \unit{\metre}).
\begin{equation}\label{eq:AppendixEulerCriticalLoad}
    P_{cr}=\frac{\pi EI}{(KL)^2}
\end{equation}
Young's moduli for conventional thermoset polyesters were expected to range
 from 2-4.5~\unit{\giga\pascal} \autocite{Lawrence1960}.
The second moment of the area for the top surface of the cylindrical sample
 with radius 6.35~\unit{\milli\metre} was calculated to be
 \(1.276\times10^{-9}\)~\unit{\metre^{4}} by assuming it to be a perfect
 filled circular area, for which the generic equation
 \ref{eq:AppendixCircularSecondMomentOfArea} applied (where the \(r\) is the
 radius of the cylinder in \unit{\metre}).
\begin{equation}\label{eq:AppendixCircularSecondMomentOfArea}
    I_{y}=\frac{\pi}{4}r^{4}
\end{equation}
By assuming that both ends of the column would be braced against rotation
 or translation (given their compression between two platens), the
 theoretical column effective length factor \(K\) would have been 0.5; as a
 further precaution a conservative design value of 0.65 was assumed
 based on pertinent literature \autocite{AISC1989}.
On these assumptions it was estimated that the samples in question would
 be likely to undergo Euler buckling at an applied load of between 
 433-976~\unit{\kilo\newton}.
As the compressive strength expected of conventional unsaturated
 polyester thermosets are expected to fall within the range of
 89-252~\unit{\mega\pascal}, corresponding to maximum compressive
 forces of between 11-32~\unit{\kilo\newton} (given the surface area of
 the samples), we concluded that the samples were unlikely to
 experience Euler buckling.

\paragraph{Correction of Raw Data}
Finite compliance of the testing machine introduced an additional 
 displacement contribution into raw exported datasets.
To correct for this, a blank compression test was performed with no sample 
 positioned between the platens, generating a force-displacement curve 
 corresponding to the inherent compliance of the machine.
This curve was subsequently used as a correctional dataset.
Equation \ref{eq:AppendixCorrectedDisplacement} shows how the corrected displacement
 value (\(d\) in \unit{\milli\metre}) was obtained by taking a
 raw displacement value (\(d_{raw}\) in \unit{\milli\metre}) and
 subtracting the corresponding correctional displacement value
 (\(d_{cor}\)).
\begin{equation}\label{eq:AppendixCorrectedDisplacement}
    d = d_{\mathrm{raw}} - d_{\mathrm{cor}}
\end{equation}\

\paragraph{Processing of Corrected Data}
Load values from a sample's corrected dataset were converted into
 engineering stress values by dividing by the initial surface area.
Each cylindrical sample's diameter was measured using a ``Magnusson''
 digital caliper (accuracy \(\pm\)~0.02~\unit{\milli\metre}) from three
 different positions (\(\text{ø}_{1},\text{ø}_{2},\text{ø}_{3}\)
 in \unit{\milli\metre}), each offset by roughly +45\degree~around the
 circumference, and the mean of these values (\(\text{ø}_{\mathrm{avg}}\) in
 \unit{\milli\metre}) calculated (see Equation
 \ref{eq:AppendixAverageTopSurfaceDiameter}).
\begin{equation}\label{eq:AppendixAverageTopSurfaceDiameter}
    \text{ø}_{\mathrm{avg}} = \frac{\text{ø}_{1}+\text{ø}_{2}+\text{ø}_{3}}{3}
\end{equation}
Equation \ref{eq:AppendixAverageTopSurfaceArea} shows how the initial top-face
 surface area was calculated from the average top surface diameter
 \(\text{ø}_{\mathrm{avg}}\).
 \begin{equation}\label{eq:AppendixAverageTopSurfaceArea}
    A_{\mathrm{init}} = 0.25 \times \pi \times \text{ø}_{\mathrm{avg}}^{2}
\end{equation}
Subsequently, the engineering stress (\(\sigma\) in \unit{\mega\pascal}) could
 be calculated (Equation \ref{eq:AppendixEngineeringStress}) by considering the
 force applied (\(F\) in \unit{\newton}) and the initial top surface area,
 \(A_{\mathrm{init}}\) \autocite{Martin2006}.
\begin{equation}\label{eq:AppendixEngineeringStress}
    \sigma _{\mathrm{eng}} =\frac{F}{A_{\mathrm{init}}}
\end{equation}

Displacement values from a sample's corrected dataset were converted into
 engineering strain by taking the difference between the initial length of 
 the samples and the length of the sample during compression.
The initial length (\(l_{\mathrm{init}}\) in \unit{\milli\metre}) of sample 
 was calculated by taking the mean of three length measurements 
 (\(l_{1},l_{2},l_{3}\) in \unit{\milli\metre}) (these taken in the same 
 fashion as for diameter) as shown in Equation~\ref{eq:AppendixAverageInitialSampleLength}.
\begin{equation}\label{eq:AppendixAverageInitialSampleLength}
    l_{\mathrm{init}} = \frac{l_{1}+l_{2}+l_{3}}{3}
\end{equation}
Equation~\ref{eq:AppendixSampleLengthUnderCompression} shows the simple link between
 length of a sample at a given point during a compressive test (\(l\) in
 \unit{\milli\metre}), the initial sample length (\(l_{\mathrm{init}}\)) and the
 corrected displacement at that point \(d\).
\begin{equation}\label{eq:AppendixSampleLengthUnderCompression}
    l = l_{\mathrm{init}}-d
\end{equation}
Subsequently, the engineering strain was calculated by finding the 
 absolute change in sample length (\(\Delta l_{\mathrm{init}}\) in 
 \unit{\milli\metre}) and dividing this by the initial length of the
 sample (\(l_{\mathrm{init}}\)) as shown in 
 Equation~\ref{eq:AppendixEngineeringStrain}.
\begin{equation}\label{eq:AppendixEngineeringStrain}
    \epsilon _{\mathrm{eng}}=\frac{\Delta l_{\mathrm{init}}}{l_{\mathrm{init}}}=\frac{l-l_{\mathrm{init}}}{l_{\mathrm{init}}}
\end{equation}

\paragraph{Stress-Strain Smoothing}
To produce a smoothed dataset (from which Young's modulus and yield 
 strength could be deduced) a combination of sample-rate reduction and 
 filtering were applied to downsample the stress-strain dataset.
Downsampling was achieved using a windowed averaging technique, whereby
 the unsmoothed dataset (\(p_{\mathrm{us}}\)) was
 split  into a number of windows along
 the strain-axis (\(p_{\mathrm{s}}\)) by electing a suitable number of points per
 window (\(p_{\mathrm{w}}\)), as shown in Equation~\ref{eq:AppendixWindowDivisionOfDataset}.
For convenience, \(p_{\mathrm{us}}\) is required to be divisible by
 \(p_{\mathrm{w}}\), and \(p_{\mathrm{s}}> 200\) was found to give reliable
 results.
Given the constant size of the unsmoothed dataset (\(p_{\mathrm{us}}=1,000\)),
 a constant number of points per window was elected (\(p_{\mathrm{w}}=5\)).
\begin{equation}\label{eq:AppendixWindowDivisionOfDataset}
    p_{\mathrm{s}}=p_{\mathrm{us}}/p_{\mathrm{w}}
\end{equation}
Equation~\ref{eq:AppendixWindowAveraging} shows how the processs of averaging stress
 (\(\sigma_{\mathrm{eng}}^{i}\)) and strain (\(\epsilon_{\mathrm{eng}}^{i}\)) values within
 a given window yielded a single point described by its average stress
 (\(\overline{\sigma}_{\mathrm{eng}}^{j}\)) and strain (\(\overline{\epsilon}_{\mathrm{eng}}^{j}\)).
\begin{equation}\label{eq:AppendixWindowAveraging}
    \overline{\sigma}_{\mathrm{eng}}^{j},\overline{\epsilon}_{\mathrm{eng}}^{j}=\frac{\sum_{i=1}^{p_{\mathrm{w}}}\sigma_{\mathrm{eng}}^{i}}{p_{\mathrm{w}}},\frac{\sum_{i=1}^{p_{\mathrm{w}}}\epsilon_{\mathrm{eng}}^{i}}{p_{\mathrm{w}}}
\end{equation}
Averaged points from each window were assembled to form the new downsampled
 dataset depicted in Equation~\ref{eq:AppendixSmoothedDataset}.
\footnotesize
\begin{equation}\label{eq:AppendixSmoothedDataset}
    [[\overline{\sigma}_{\mathrm{eng}}^{j},\overline{\epsilon}_{\mathrm{eng}}^{j}],[\overline{\sigma}_{\mathrm{eng}}^{j+1},\overline{\epsilon}_{\mathrm{eng}}^{j+1}],...,[\overline{\sigma}_{\mathrm{eng}}^{p_{\mathrm{s}}},\overline{\epsilon}_{\mathrm{eng}}^{p_{\mathrm{s}}}]]
\end{equation}
\normalsize
Figure~\ref{fig:AppendixAverageSmoothing} demonstrates the efficacy of this
 technique in removing noise from the unsmoothed stress-strain dataset.
\noindent
\begin{minipage}{0.485\textwidth}
    \includegraphics[width=\textwidth]{AverageSmoothing.png}
    \captionsetup{format=custom,width=\textwidth,skip=0pt}
    \captionof{figure}{
        \footnotesize{
            Multi-scale depiction of the stress-strain dataset before
            and after downscaling.
            (A) Overview line plot.
            (B) Zoomed scatter plot.
        }
    }
    \label{fig:AppendixAverageSmoothing}
\end{minipage}

\paragraph{Derivative Analysis}
The proportional limit is defined as the stress at which a material
 ceases to obey Hooke's law (i.e. stress no longer increases proportionally
 with strain).
The gradient of the preceding elastic region of deformation
 is generally defined as the the Young's modulus.
However, given that the stress-strain curves did not consistently exhibit
 a linear region, the limit of proportionality was equated to the
 point at which the stress-strain curve began to level off (i.e. the first
 point of inflection).
To find this inflection point, the first derivative of stress with respect
 to strain (i.e. tangent modulus \(E_{\mathrm{t}}\) in \unit{\mega\pascal}) was calculated
 from the smoothed stress-strain dataset by
 taking into account the instantaneous change in both stress (\(\Delta
 \sigma_{\mathrm{eng}}\)) and strain (\(\Delta \epsilon_{\mathrm{eng}}\)), and
 given by Equation~\ref{eq:AppendixTangentModulus}.

\begin{equation}\label{eq:AppendixTangentModulus}
    E_{\mathrm{t}}=\frac{\Delta \sigma_{\mathrm{eng}}}{\Delta \epsilon_{\mathrm{eng}}}=\frac{\sigma_{\mathrm{eng}}^{j+1}-\sigma_{\mathrm{eng}}^{j}}{\epsilon_{\mathrm{eng}}^{j+1}-\epsilon_{\mathrm{eng}}^{j}}
\end{equation}

Figure~\ref{fig:AppendixDerivativeComparison} shows how smoothing of the 
 stress-strain dataset facilitated the analysis of the resultant tangent
 modulus versus strain dataset.
Step-like features are apparent in the tangent-modulus curve obtained from
 the unsmoothed stress-strain data.
These may be consistent with either distretisation or internal processing
 of the exported displacement signals and subsequently amplified by
 the process of numerical differentiation.
To avoid undue physical significance being assigned to these artefacts,
 proportional-limit detection was performed upon a smoothed stress-strain
 dataset.

\noindent
\begin{minipage}{0.485\textwidth}
    \includegraphics[width=\textwidth]{DerivativeComparison.png}
    \captionsetup{format=custom,width=\textwidth,skip=0pt}
    \captionof{figure}{
        \footnotesize{
            A comparison of the tangent modulus vs strain curves obtained
            from both the raw and the smoothed counterpart stress-strain
            datasets.
            (A) Strain at first inflection point.
        }
    }
    \label{fig:AppendixDerivativeComparison}
\end{minipage}

Detection of the first inflection point was automated using an
 algorithm that iterated chronologically through the tangent modulus versus
 strain dataset from low to high strains.
At each point, the algorithm calculates an average gradient exhibited
 by an equal number of points both before and after the current strain
 value (the `point horizon').
Using a point horizon of ten, the first inflection point was defined as
 the first point at which the summed forward and backward gradients fell
 below zero.
This point was taken as the proportional limit, as shown in 
 Figure~\ref{fig:AppendixDerivativeComparison}A.

\subsection{Notes on Gaussian Processes}
Figure~\ref{fig:AppendixSliceGP} presents a matrix of slice plots involving the two
 Gaussian processes as they existed at the culmination of the multi-objective
 experiment.
Every plot holds \nicefrac{2}{3} predictors constant whilst the column
 dictates which of the three predictors is varied.
Each row of the matrix conditions the Gaussian process \(\mathcal{GP}\) on a
 different set of predictor constants, these are derived from either the 
 14\textsuperscript{th} sample (\(\mathbf{x}^{i=14}\)), the
 68\textsuperscript{th} sample (\(\mathbf{x}^{i=68}\)), or the
 mean of all samples' predictors (\(\mathbf{x}^{\mu}\)).
The graphs are further annotated with vertical lines representing sample
 predictor values (\(\{x_{1},x_{2},x_{3}\}\)):
 \(\mathbf{x}^{i=14}=\{0.98,0.5,0.93\}\)
 \(\mathbf{x}^{i=68}=\{0.72,0.4,0.5\}\)
 \(\mathbf{x}^{\mu}=\{0.71,0.64,0.64\}\).


Row B of Figure~\ref{fig:AppendixSliceGP} describes the processes as conditioned
 upon the experimental mean of sample predictors
 (\(\mathcal{GP}|\mathbf{x}^{\mu}\)); these graphs provide rough
 intuition with respect to the objective topography of the parameter space
 but naturally fail to account for why specific combinations of
 predictor values (e.g. sample 68 or 14) exhibit higher or lower performance
 in both the metrics of interest.

In Figure~\ref{fig:AppendixSliceGP}A we see the processes conditioned on the
 68\textsuperscript{th} sample's predictors
 \(\mathcal{GP}|\mathbf{x}^{i=68}\): in the case of predictor \(x_{3}\)'s
 yield strength process (Fig.\ref{fig:AppendixSliceGP}:A3) we see a topographical
 deviation from the mean conditioned process (Fig.\ref{fig:AppendixSliceGP}:B3).
This deviation was a flattening in the region where \(x_{3}=[0.5,0.65]\)
 which shifted the local optimum toward \(x_{3}=0.5\): this resulted in an
 alignment of both local optima of the two processes, which partly explains
 why marginalisation of UPR content by sample 68 leads to improved
 performance in both objective metrics.

Although the troughs apparent in the mean processes with respect to \(x_{1}\)
 and \(x_3\) (Fig.\ref{fig:AppendixSliceGP}:B1\&B2) hinted at why the
 14\textsuperscript{th} sample's predictor value combination produced
 suboptimal performance in both the objective metrics, the \(x_{2}\)
 topography failed to do so.
On examination of the \(x_{2}\) topography of the processes conditioned on
 the 14\textsuperscript{th} sample's predictors
 \(\mathcal{GP}|\mathbf{x}^{i=14}\) we find large deviations from the mean
 processes: both the topographies have fallen noticably, the Young's modulus
 process has flattened, and the yield strength process inverted such that
 the local optimum shifted from \(x_{3}=0.4\) to \(x_{3}=1\).
Furthemore, this inversion presented a local trough in the vicinity of the
 14\textsuperscript{th} sample explaining further why this sample behaved
 suboptimally.

\begin{figure*}[htb]
    \captionsetup{format=custom,width=0.8\textwidth,skip=0pt}
    \includegraphics[width=\textwidth]{MOBO-ConditionalMatrix.png}
    \captionof{figure}{
            Slice plot of the Gaussian processes developed by the
            multi-objective experiment with mean function and a
            95\% confidence interval (\(1.96\sigma\)) presented.
            UP, Unsaturated polyester.
            RD, Reactive diluent.
            DBP, Dibenzoyl peroxide.
            CS, Calcium sulphate.
            UPR, Unsaturated polyester resin.
            IA, Initiating agent.
    }
    \label{fig:AppendixSliceGP}
\end{figure*}

These three observations are respectively reflected in the Gaussian
 process topographies presented in Figure~\ref{fig:SliceGP}: when considering
 predictor \(x_{3}\) (column 3: controlling balance of unsaturated polyester
 resin to initiating agent) we see that Gaussian processes clearly describe
 a negative correlation between \(x_{3}\) and optimality in all three slice
 plots (A3,B3,C3).
These strong correlations and coinciding local optima
 between \(\mathcal{GP}\)'s would explain why there is a clear global link
 between the balance of unsaturated polyester resin and initiating agent and
 sample optimality,
With respect to \(x_{2}\) (column 2: controlling initiating agent composition)
 the \(\mathcal{GP}\)'s display a negative correlation between \(x_{2}\) and
 optimality in two of the slice plots (A2,B2), but this relation disappears
 and is replaced with a flattened and lowered process in at least one region
 of the parameter space (C2).
In the case of \(x_{1}\) (column 1: controlling unsaturated polyester resin
 composition) we can see that the \(\mathcal{GP}\)'s display a function shape 
 with singular peaks at varying positions in all the different slice plots
 (A1,B1,C1); this added optimisation complexity explains well why there is
 not obvious global trend relating the unsaturated polyester resin content
 with sample optimality.